\section{Introduction}
\label{sec:intro}

Much of experimental biology is a search for interventions that move a cell from one state to
another: differentiating a stem cell, repolarizing a T cell, reverting a diseased phenotype,
or---in drug discovery---combining perturbations to reprogram a pathogenic program. When single
perturbations are insufficient, the field turns to \emph{combinations}, and the design space
explodes super-linearly: a library of $100$ single perturbations already implies roughly
$5{,}000$ pairs and $160{,}000$ triples. No screen can measure them all, and most combinations
are not worth measuring---their joint effect is what the single effects already predict. The
decisive questions are therefore triage questions: \emph{which unmeasured combinations are worth
running}, and, once a combination is measured, \emph{is its effect genuinely emergent or merely
the additive sum of its parts}. These questions are rarely asked with a checkable answer, and
their economic weight is large: only about one in ten clinical programs succeeds, and many
late-stage failures involve insufficient efficacy, which motivates grounding combination choice
in measured perturbation data rather than inferred models
\citep{Minikel2024,Nelson2015,Harrison2016,BegleyEllis2012,Prinz2011}.

The computational tools available fall into two camps that, as we document in
Section~\ref{sec:related}, do not overlap. Perturbation-response models---from
gene-regulatory-network simulators to single-cell foundation models---are increasingly trained
on real screens but are forward predictors: given a combination they return a predicted response,
and a predictor structurally always returns \emph{some} answer, so it cannot certify that a
combination reaches a state the singles cannot. Network-control and Boolean-attractor methods do
reason about reachability and can carry controllability guarantees, but those guarantees are
properties of a hand-built or inferred model, only as trustworthy as the network assumed. Neither
camp certifies emergence from measured effects. This gap matters all the more because a wave of
2025--2026 benchmarks finds that current perturbation predictors, foundation models included,
frequently fail to beat mean or linear baselines on genuinely unseen perturbations
\citep{AhlmannEltze2025,Wong2025,Kernfeld2025,Atti2025}. We reproduce this directly: a
self-contained learned neural baseline ties the trivial additive baseline for blind prediction of
held-out doubles (cosine $0.896$ vs.\ $0.897$; Section~\ref{sec:results-headtohead}). Accuracy,
in other words, is the wrong axis on which to claim a combination is special---the field's own
evidence argues for an output that is weaker to state but falsifiable.

\paragraph{The convex-cone move.}
We take the opposite starting point from model building: we do not infer a network at all. A
combinatorial perturbation screen directly measures, for each single perturbation, the effect
vector it produces in expression space. We treat these single-gene vectors as a fixed dictionary
and ask a geometric question of each combination: does its measured effect lie in the set of
shifts the singles can produce? Two modelling choices make this a convex-geometry problem. A
measured single effect may be applied with non-negative weight but not reversed into an unmeasured
``negative perturbation''; and several perturbations applied together are approximated, to first
order, by the sum of their single effects. Individual effect vectors may raise some transcripts
and lower others. The set of directions reachable by the singles is therefore the
\emph{non-negative conic hull} of the measured single-gene effect vectors, and asking whether a
combination is emergent is asking whether its measured effect lies \emph{outside} that
model-relative cone---a question non-negative least squares (NNLS) answers exactly
\citep{LawsonHanson1974}. This reframing turns ``is this combination special?'' into a decision
with a checkable answer on both sides. When the combination lies inside the cone, its effect is
reachable by the singles and the NNLS weights name the reachable mixture. When it lies outside,
convex duality supplies a separating hyperplane---a Farkas/KKT certificate
\citep{Farkas1902,KuhnTucker1951}---whose full vector certifies that no non-negative mixture of
the singles reaches that direction.

\paragraph{The honest bar.}
Geometric ``outside'' is necessary but not sufficient. Because the unreachable residual is a
normalized quantity, a low-magnitude combination measured with the same absolute noise as a
large one shows an inflated apparent residual: on the Norman screen the raw unreachable fraction
correlates with effect magnitude at Spearman $-0.56$, and only a minority of doubles clear their
own split-half noise floor. We therefore make the certificate noise-aware: a per-gene
noise-injection null asks whether a reachable target, corrupted by the measured noise, would look
as far outside the cone as the observed combination does. This removes the magnitude confound
(the noise-aware score correlates with magnitude at $+0.14$) and yields a two-bar verdict---a
permissive significance bar and a stricter effect-size bar---that a learned predictor, emitting a
point estimate with no null, cannot supply.

\paragraph{Scope.}
The method is a decision instrument, not an oracle of biological truth: it certifies emergence
\emph{relative to the measured single-gene dictionary} under an explicit additivity model, and
its verdicts are hypotheses for the bench, not proofs of biological synergy. The synergistic
pairs it surfaces (e.g.\ the CEBPA/CEBPE differentiation program) are already reported by the
source screen \citep{Norman2019} and enter here as positive controls; the novelty we claim is the
measured-effect, certificate-carrying \emph{triage decision} itself. We report the source data's
provenance verbatim: the Norman combinatorial CRISPRa file used here is labelled A549, whereas the
canonical Norman 2019 line is K562, and we carry that discrepancy through every result.

\paragraph{Contributions.}
This paper makes four contributions.
\begin{enumerate}
\item \textbf{A certified emergence test for combinatorial screens.} We formalize emergence of a
combination as \emph{non-}membership in the convex cone of measured single-gene effect vectors,
and solve it with NNLS so that a verdict and---on emergence---a Farkas/KKT dual certificate follow
from one convex program, made noise-aware by an injection null that removes the magnitude confound
(Section~\ref{sec:methods}).
\item \textbf{A validated headline result.} On the Norman combinatorial CRISPRa screen, all $131$
measured doubles are geometrically outside the single-gene cone; the noise-aware null certifies
$129/131$ at $p<0.05$ ($35$ also clearing a $1.9\times$ floor) and recovers classical synergy at
partial Spearman $0.62$ controlling for magnitude (Section~\ref{sec:results-certificate}).
\item \textbf{Training-free prospective triage, and why accuracy is the wrong axis.} A single
scalar computed from the two single effects---their negative cosine---ranks emergent combinations
$2.4\times$ above base rate without measuring them, while a self-contained learned MLP ties the
additive baseline for blind prediction yet emits no certificate on any of the $131$ outside-cone
doubles (Sections~\ref{sec:results-triage}--\ref{sec:results-headtohead}).
\item \textbf{Generality across order and across screens.} The machinery extends to arbitrary
interaction order $k$---a reachable-from-lower-order test shrinks the certified doubles $40\to16$
on real data, and a synthetic planted-epistasis triple screen recovers order-3 emergence at AUC
$1.00$---and the certificate with its noise-aware discipline transfers to an orthogonal screen
(CaRPool-seq, Cas13d knockdown in THP-1), where the raw-label triage rank also transfers but the
magnitude-controlled component is screen-specific
(Sections~\ref{sec:results-kway}--\ref{sec:results-crossscreen}).
\end{enumerate}
